\section*{Введение}
\addcontentsline{toc}{section}{Введение}

В этой главе основная теорема теории Галуа (\rf{3.17}) применяется к многочленам и расширениям полей. Почти каждый раздел
устроен одинаково: вопрос о полях (существует ли элемент, корень, базис с нужными свойствами) переводится в вопрос о
конечной группе Галуа и её подгруппах, а там решается методами теории групп или линейной алгебры.

\begin{itemize}[leftmargin=*]
\item \emph{Примитивный элемент} (\rf{5.1}--\rf{5.5}). Каждое конечное сепарабельное расширение порождается одним элементом;
простота конечного расширения равносильна конечности множества промежуточных полей. Пример $k(X^p,Y^p)\subset k(X,Y)$
показывает, что без сепарабельности утверждение неверно.
\item \emph{Основная теорема алгебры} (\rf{5.6}--\rf{5.7}): поле $\CC$ алгебраически замкнуто. Из анализа нужны лишь два свойства
$\R$, всё остальное делают теоремы Силова и соответствие Галуа.
\item \emph{Круговые расширения} (\rf{5.8}--\rf{5.13}): поля $\Q(\zeta_n)$, круговые многочлены $\Phi_n$ и их неприводимость,
изоморфизм $\Gal(\Q(\zeta_n)/\Q)\cong(\Z/n\Z)^\times$, подполя круговых полей, суммы Гаусса, элементы Фробениуса (и с ними
квадратичный закон взаимности) и теорема Гаусса о построении правильных многоугольников.
\item \emph{Независимость характеров} Дедекинда (\rf{5.14}--\rf{5.16}) и \emph{теорема о нормальном базисе} (\rf{5.17}--\rf{5.19}):
конечное расширение Галуа $E/F$ с группой $G$ изоморфно групповой алгебре $F[G]$ как $F[G]$-модуль.
\item \emph{$G$-модули и когомологии} (\rf{5.20}--\rf{5.26}): скрещённые гомоморфизмы, группа $H^1(G,M)$ и \emph{теорема
Гильберта~90} $H^1(G,E^\times)=0$; для циклического расширения она означает, что элементы нормы $1$~--- это в точности
элементы вида $\beta/\sigma\beta$, где $\sigma$~--- образующая группы Галуа.
\item \emph{Циклические расширения} (\rf{5.27}--\rf{5.29}) и \emph{теория Куммера} (\rf{5.30}--\rf{5.32}): если $F$ содержит первообразный
корень $n$-й степени из $1$, то конечные абелевы расширения $F$, экспонента которых делит $n$, получаются присоединением
корней $\sqrt[n]{a}$ и описываются конечными подгруппами группы $F^\times/F^{\times n}$.
\item \emph{Теория Артина--Шрайера и векторы Витта} (\rf{5.32a}--\rf{5.32x}): аналог теории Куммера в характеристике $p$ для
абелевых расширений экспоненты $p$ (корни $\sqrt[p]{a}$ заменяются корнями многочленов $X^p-X-a$) и экспоненты, делящей
$p^n$ (вместо элементов поля~--- векторы Витта длины $n$). Попутно выясняется, что уравнение $X^p-X-t=0$ над $\F_p(t)$ с
группой Галуа $\Z/p\Z$ не разрешимо в радикалах.
\item \emph{Разрешимость в радикалах} (\rf{5.33}--\rf{5.35e}): теорема Галуа~--- над полем характеристики $0$ многочлен разрешим в
радикалах тогда и только тогда, когда его группа Галуа разрешима; в характеристике $p$ к радикалам нужно добавить корни
многочленов $X^p-X-a$ (\rf{5.35b}). Примеры~--- уравнения степеней $3$, $4$ и $5$.
\item \emph{Симметрические многочлены} (\rf{5.36}--\rf{5.39}) и \emph{общий многочлен степени $n$} (\rf{5.40}--\rf{5.42}): группа Галуа общего
многочлена~--- $S_n$, поэтому в характеристике $0$ при $n\ge5$ формулы в радикалах для его корней нет.
\item \emph{Норма и след} (\rf{5.43}--\rf{5.50}), включая несепарабельный случай.
\end{itemize}
